\documentclass[11pt,letterpaper]{report}

% CSC / CTC 492 - Senior Design
% Working research-paper skeleton. This draft contains planned sections,
% not completed experiments or claimed results.
\usepackage[margin=1in]{geometry}
\usepackage[T1]{fontenc}
\usepackage[utf8]{inputenc}
\usepackage{lmodern}
\usepackage{microtype}
\usepackage{graphicx}
\usepackage{booktabs}
\usepackage{listings}
\usepackage{xcolor}
\usepackage{hyperref}
\usepackage[backend=biber,style=numeric-comp,sorting=none]{biblatex}
\addbibresource{references.bib}

\hypersetup{
    colorlinks=true,
    linkcolor=blue!45!black,
    citecolor=blue!45!black,
    urlcolor=blue!45!black,
    pdftitle={Document Extractor: Senior Design Research Paper}
}

\lstset{
  basicstyle=\ttfamily\small,
  breaklines=true,
  frame=single,
  numbers=left,
  numberstyle=\tiny,
  showstringspaces=false
}

\newcommand{\todowrite}[1]{\par\noindent\textit{Draft placeholder: #1}\par}

\begin{document}
\pagenumbering{gobble} % prevent duplicate PDF page anchors from title matter

% Cover page - update the placeholders below before submitting.
\begin{titlepage}
\centering
\vspace*{1.3in}
{\Large California State University, Dominguez Hills\par}
\vspace{0.3in}
{\large CSC / CTC 492 -- Senior Design\par}
\vspace{0.9in}
{\Huge\bfseries Document Extractor\par}
\vspace{0.25in}
{\Large Structured Document Extraction, Validation, and Consistency Checking\par}
\vspace{1.1in}
{\large Student: [Your Name]\par}
\vspace{0.15in}
{\large Instructor: [Confirm Instructor Name]\par}
\vfill
{\large Fall 2026 --- Working Draft\par}
\end{titlepage}

\pagenumbering{roman}
\chapter*{Abstract}
\addcontentsline{toc}{chapter}{Abstract}
This project studies the extraction and validation of structured information from
unstructured documents. An initial Python prototype processes a predefined set
of fields from sample intake documents and applies field-validation rules.
The planned extension adds work-order and invoice processing to investigate
whether related documents can be compared to identify potentially missing or
inconsistent entries. The system design and evaluation will emphasize traceable
findings, human review, and measurable detection performance.
\todowrite{Revise this abstract after experiments are complete. Include the final method, actual measured results, limitations, and contributions.}

\tableofcontents
\clearpage
\pagenumbering{arabic}

\chapter{Introduction}
\section{Motivation}
Organizations use documents such as intake forms, work orders, and invoices
to record business information. Transcribing and checking these documents can
require manual review. This project examines how structured extraction and
validation can support a more consistent review process.
\todowrite{Add a documented motivation with sources and define the actual target workflow.}

\section{Research Question}
\emph{How reliably can a document-processing system extract fields from
work orders and invoices and flag possible inconsistencies between related
records for human review?}
\todowrite{Refine the research question before collecting experimental data.}

\section{Scope and Planned Contributions}
The current prototype centers on extracting fields and applying validation
rules. Planned work extends it to additional document schemas and
cross-document consistency checks. No performance advantage is claimed until
it is measured.

\chapter{Background and Related Work}
\section{Document Information Extraction}
\todowrite{Review relevant peer-reviewed work on document extraction, layout analysis, and structured field prediction. Cite sources actually read.}

\section{Validation and Human Review}
\todowrite{Review methods for field validation, uncertainty, human-in-the-loop correction, and traceability.}

\section{Cross-Document Consistency Checking}
\todowrite{Compare existing approaches for aligning records, entity matching, and discrepancy detection.}

\chapter{Method}
\section{System Overview}
The existing prototype is implemented in Python. It defines expected fields,
extracts structured values, and applies validation rules. The research
extension will introduce representations for work orders and invoices and
compare aligned records.
\todowrite{Insert the actual architecture diagram and explain how data moves between components.}

\section{Document Inputs and Data Schemas}
\todowrite{Specify accepted input types, schema fields, normalization rules, and representative examples. Distinguish implemented from planned schemas.}

\section{Extraction and Validation Pipeline}
\todowrite{Describe the actual extraction steps, model/API configuration, validation rules, uncertainty handling, and error cases.}

\section{Proposed Consistency-Check Module}
The proposed first check compares documented materials or services against
corresponding invoice entries. A missing or mismatched item is treated as a
\emph{potential discrepancy}, not as confirmed unbilled work, until reviewed.
\todowrite{Define record matching, tolerances, deterministic checks, ambiguous cases, and evidence references after implementation.}

\section{Reproducibility and Implementation Environment}
\todowrite{Record exact software versions, commands, datasets, experiment parameters, and how to reproduce every measurement.}

\chapter{Results}
\section{Evaluation Dataset}
\todowrite{Describe the dataset, source/permission status, number of examples, ground-truth labels, and train/test separation (if applicable).}

\section{Evaluation Metrics}
Planned measurements include the number of true discrepancy findings,
false-positive alerts, missed discrepancies, and processing time. The final
metric definitions and evaluation protocol will be fixed before results are
reported.
\todowrite{Define precision, recall, false-positive rate, and timing procedure for the actual experiment.}

\section{Experimental Findings}
\todowrite{Insert only measured tables/figures with reproducible source data. No experiments have been reported in this starter draft.}

\chapter{Discussion}
\section{Interpretation}
\todowrite{Explain what measured results mean after the experiments are complete.}

\section{Limitations and Threats to Validity}
\todowrite{Discuss input quality, synthetic-vs-real data, uncertain matches, potential false alerts, and generalization limits.}

\chapter{Conclusion and Future Work}
\todowrite{Summarize supported findings, not planned outcomes. Describe future document schemas, review workflows, and integrations separately from completed work.}

% Add verified entries to references.bib and cite them with \cite{key}.
% Remove the sentence below once bibliography entries have been added.
\chapter*{References}
\addcontentsline{toc}{chapter}{References}
\noindent\textit{Bibliographic references will be added after sources are read and verified.}
\printbibliography[heading=none]

\chapter*{Glossary}
\addcontentsline{toc}{chapter}{Glossary}
\begin{description}
  \item[Document extraction] Converting relevant information from a document into structured fields.
  \item[Schema] A definition of expected fields and data types.
  \item[Validation rule] A test that checks data against an explicitly defined constraint.
  \item[Discrepancy] A possible inconsistency between related records requiring assessment.
  \item[Human review] A person checks a flagged result before it is treated as confirmed.
\end{description}

\appendix
\chapter{Supplementary Material}
\todowrite{Add sample input descriptions, additional tables, reproducibility details, and experiment outputs as needed.}

\end{document}
